    \documentclass[conference]{IEEEtran}
    
    % ---------------------------------------------------------
    % PACKAGES
    % ---------------------------------------------------------
    
    \usepackage{cite}
    \usepackage{amsmath,amssymb,amsfonts}
    \usepackage{graphicx}
    \usepackage{booktabs}
    \usepackage{multirow}
    \usepackage{array}
    \usepackage{url}
    \usepackage{xcolor}
    \usepackage{microtype}
    
    % ---------------------------------------------------------
    % DOCUMENT INFORMATION
    % ---------------------------------------------------------
    
    \title{
    Switch FI: Evaluating Memory-Augmented Grounded Reasoning
    for Personalized Financial LLM Assistance
    }
    
    \author{
    \IEEEauthorblockN{
    Author Name\IEEEauthorrefmark{1},
    Author Name\IEEEauthorrefmark{2},
    Author Name\IEEEauthorrefmark{3}
    }
    \IEEEauthorblockA{
    \IEEEauthorrefmark{1}Department / Institution\\
    City, Country\\
    email@example.com
    }
    \and
    \IEEEauthorblockN{
    Author Name\IEEEauthorrefmark{4}
    }
    \IEEEauthorblockA{
    \IEEEauthorrefmark{4}Department / Institution\\
    City, Country\\
    email@example.com
    }
    }
    
    % ---------------------------------------------------------
    % BEGIN DOCUMENT
    % ---------------------------------------------------------
    
    \begin{document}
    
    \maketitle
    
    % ---------------------------------------------------------
    % ABSTRACT
    % ---------------------------------------------------------
    
    \begin{abstract}
    
    Large language models (LLMs) can provide useful financial guidance but
    remain unreliable when answering questions requiring access to a user's
    actual financial records or persistent knowledge of user-specific
    preferences. In particular, ungrounded LLM responses can hallucinate
    numerical values, while stateless assistants repeatedly lose information
    about user preferences and financial goals.
    
    We present \textit{Switch FI}, a Model Context Protocol (MCP)-based
    financial assistance system that combines deterministic financial
    aggregation with explicit and implicit user memory. The system connects
    LLM assistants to consented financial data through India's
    Account Aggregator (AA) ecosystem and exposes structured financial
    operations as callable tools. Rather than allowing the language model to
    directly infer numerical financial state, Switch FI performs deterministic
    aggregation over transaction data and supplies the resulting values to
    the
    model as grounded context.
    
    We further introduce an evaluation framework for measuring the effect of
    persistent memory on personalized financial reasoning. The evaluation
    compares memory-enabled and stateless configurations over a fixed set of
    financial questions spanning numerical grounding, preference adherence,
    and multi-turn memory retrieval. The experimental dataset contains
    synthetic but realistic financial records with known ground truth,
    enabling reproducible evaluation.
    
    \textbf{Results will be populated after completion of the experimental
    evaluation.}
    
    \end{abstract}
    
    \begin{IEEEkeywords}
    large language models, financial AI, personalization, memory,
    grounded reasoning, Model Context Protocol, account aggregator,
    financial question answering
    \end{IEEEkeywords}
    
    % =========================================================
    % 1. INTRODUCTION
    % =========================================================
    
    \section{Introduction}
    
    Large language models (LLMs) have demonstrated strong capabilities in
    natural language understanding, reasoning, and interactive assistance.
    However, their usefulness in personal finance is constrained by two
    fundamental limitations: lack of direct access to current financial
    records and limited persistence of user-specific information.
    
    A user may ask questions such as ``How much did I spend on food last
    month?'', ``Am I spending more than usual on subscriptions?'', or
    ``Given my preference to keep an emergency fund of six months, how much
    should I save this month?''. These questions require more than generic
    financial knowledge. They require access to user-specific transaction
    data, deterministic numerical computation, and knowledge of preferences
    that may have been expressed during earlier interactions.
    
    Without external grounding, an LLM may produce plausible but incorrect
    numerical values. Similarly, without persistent memory, an assistant may
    know that a user previously expressed a financial preference but fail to
    apply that preference in a later interaction.
    
    Recent work on LLM personalization has demonstrated the importance of
    memory extraction, retrieval, and utilization. AlpsBench, for example,
    evaluates personalization through tasks covering information extraction,
    information update, information retrieval, and information utilization
    \cite{alpsbench}. Its evaluation framework demonstrates that remembering
    information and successfully applying it are distinct capabilities.
    
    Financial assistance introduces an additional requirement: responses
    must be grounded in an external source of financial truth. This creates
    an intersection between two problems:
    
    \begin{enumerate}
        \item \textbf{Grounded financial reasoning:} obtaining accurate
        financial facts from structured financial data.
        \item \textbf{Personalized reasoning:} retrieving and applying
        user-specific preferences and constraints across interactions.
    \end{enumerate}
    
    Switch FI addresses these requirements through a hybrid architecture.
    Financial state is computed using deterministic tools, while
    personalization information is represented using structured explicit and
    implicit memory. The resulting context is then supplied to an LLM through
    an MCP interface.
    
    The central research question investigated in this work is:
    
    \begin{quote}
    \textit{Does persistent user memory improve personalized financial
    reasoning when the underlying financial facts are deterministically
    grounded?}
    \end{quote}
    
    This question is investigated using a controlled ablation between
    memory-enabled and stateless configurations.
    
    \subsection{Contributions}
    
    This work makes the following contributions:
    
    \begin{itemize}
    
        \item We present the architecture of \textit{Switch FI}, an
        LLM-agnostic MCP server for grounded personal financial assistance.
    
        \item We introduce a deterministic financial reasoning layer that
        computes financial summaries from structured transaction records
        rather than relying on LLM-generated numerical calculations.
    
        \item We implement a structured memory layer supporting both
        explicit user preferences and implicitly inferred preferences.
    
        \item We define an evaluation protocol separating numerical
        grounding, preference adherence, and multi-turn memory retrieval.
    
        \item We formulate a controlled memory ablation comparing identical
        financial data and questions under memory-enabled and stateless
        configurations.
    
    \end{itemize}
    
    % =========================================================
    % 2. RELATED WORK
    % =========================================================
    
    \section{Related Work}
    
    \subsection{LLM Personalization and Memory}
    
    Personalization in LLM systems has increasingly shifted from static
    prompt-based personalization toward persistent memory mechanisms.
    Recent benchmarks distinguish between systems that simply retrieve
    historical context and systems that explicitly evaluate the lifecycle
    of personalized information.
    
    AlpsBench evaluates four stages of personalization: personalized
    information extraction, personalized information update, personalized
    information retrieval, and personalized information utilization
    \cite{alpsbench}. The benchmark further distinguishes multiple forms of
    personalization, including preference following and constraint
    following.
    
    This distinction is important for financial assistants. A system may
    correctly remember that a user prefers conservative spending but still
    fail to apply that preference when making a recommendation.
    
    Existing personalization research therefore motivates evaluating memory
    not merely as storage, but as a mechanism that influences downstream
    reasoning.
    
    \subsection{Financial LLM Grounding}
    
    LLM applications in finance introduce additional reliability concerns.
    Financial questions frequently contain numerical quantities, temporal
    constraints, and domain-specific terminology. Hallucinated values can
    therefore result in materially incorrect conclusions.
    
    A common approach to hallucination mitigation is retrieval-augmented
    generation, where external information is retrieved before generation.
    However, financial assistance requires more than document retrieval.
    Transaction-level information must often be aggregated, categorized, and
    computed before it can be meaningfully provided to an LLM.
    
    Switch FI therefore separates deterministic financial computation from
    language generation.
    
    \subsection{Consent-Based Financial Data Access}
    
    India's Account Aggregator ecosystem provides a consent-driven
    architecture for accessing financial information from participating
    financial institutions. This model enables financial information to be
    shared with a requesting entity based on user authorization.
    
    Switch FI uses the Account Aggregator model as the data-access boundary.
    For experimentation, the system uses a Finvu sandbox adapter and a
    synthetic/mock data mode.
    
    This architecture separates:
    
    \begin{enumerate}
        \item financial data acquisition,
        \item financial computation,
        \item personalization memory, and
        \item natural-language reasoning.
    \end{enumerate}
    
    \subsection{Research Gap}
    
    Existing personalization benchmarks primarily study general-purpose
    conversational interactions, while financial LLM systems often focus on
    grounding and hallucination reduction.
    
    Switch FI combines these dimensions by evaluating whether memory affects
    personalized reasoning \textit{after} financial facts have been grounded
    through deterministic computation.
    
    % =========================================================
    % 3. SWITCH FI SYSTEM
    % =========================================================
    
    \section{Switch FI System}
    
    \subsection{System Overview}
    
    Switch FI is designed as an intermediary between an LLM assistant and a
    user's financial information.
    
    The architecture consists of four major layers:
    
    \begin{enumerate}
        \item \textbf{Consent and data acquisition layer}
        \item \textbf{Deterministic financial reasoning layer}
        \item \textbf{Memory layer}
        \item \textbf{MCP interface layer}
    \end{enumerate}
    
    Figure~\ref{fig:architecture} presents the overall architecture.
    
    \begin{figure*}[t]
        \centering
        \includegraphics[
            width=0.95\textwidth
        ]{figures/system_architecture.pdf}
        \caption{
        Switch FI system architecture. Financial data is acquired through
        a consent-based adapter, transformed into deterministic financial
        summaries, combined with relevant user memory, and exposed to an
        LLM through MCP tools.
        }
        \label{fig:architecture}
    \end{figure*}
    
    \subsection{Consent and Data Acquisition}
    
    The financial data acquisition layer abstracts the underlying data
    provider through an adapter interface.
    
    For production deployment, the adapter is intended to communicate with
    the Account Aggregator ecosystem. During development and experimentation,
    a mock adapter generates deterministic synthetic transaction data.
    
    The system supports two modes:
    
    \begin{equation}
    D =
    \begin{cases}
    D_{\mathrm{mock}}, & \text{if MOCK\_MODE = true}\\
    D_{\mathrm{AA}}, & \text{otherwise}
    \end{cases}
    \end{equation}
    
    where $D_{\mathrm{mock}}$ represents synthetic financial data and
    $D_{\mathrm{AA}}$ represents data obtained through the Account Aggregator
    integration.
    
    The mock dataset is generated with known ground truth so that numerical
    answers can be evaluated without ambiguity.
    
    \subsection{Deterministic Financial Reasoning}
    
    A central design principle of Switch FI is that the LLM should not be
    responsible for computing authoritative financial totals.
    
    Let the user's transaction dataset be
    
    \begin{equation}
    T = \{t_1,t_2,\ldots,t_n\}.
    \end{equation}
    
    Each transaction $t_i$ contains fields including amount, timestamp,
    merchant, and category.
    
    A deterministic aggregation function computes a financial state:
    
    \begin{equation}
    z = f(T),
    \end{equation}
    
    where $z$ may contain:
    
    \begin{itemize}
        \item total income,
        \item total expenditure,
        \item category-wise expenditure,
        \item recurring expenses,
        \item monthly summaries,
        \item savings,
        \item account balances, and
        \item other derived financial quantities.
    \end{itemize}
    
    The resulting structured output is then supplied to the LLM.
    
    This creates the following separation:
    
    \begin{equation}
    \text{Financial Truth} \rightarrow
    \text{Deterministic Computation} \rightarrow
    \text{LLM Explanation}.
    \end{equation}
    
    For transaction categorization, the system uses a regex/rule-based
    strategy for deterministic matches, followed by an LLM fallback for
    ambiguous transactions.
    
    \subsection{Memory Architecture}
    
    Switch FI maintains structured user memory independently from raw
    financial transactions.
    
    For a user $u$, memory is represented as:
    
    \begin{equation}
    M_u = \{m_1,m_2,\ldots,m_k\}.
    \end{equation}
    
    Each memory record may contain:
    
    \begin{itemize}
        \item memory content,
        \item memory type,
        \item source,
        \item confidence,
        \item creation timestamp,
        \item update timestamp.
    \end{itemize}
    
    We distinguish two memory categories.
    
    \subsubsection{Explicit Memory}
    
    Explicit memory represents information directly stated by the user.
    
    Examples include:
    
    \begin{itemize}
        \item ``I want to maintain a six-month emergency fund.''
        \item ``I prefer not to spend more than INR 5,000 on subscriptions.''
        \item ``I am saving for a car.''
    \end{itemize}
    
    \subsubsection{Implicit Memory}
    
    Implicit memory represents preferences or behavioral information inferred
    from repeated interactions.
    
    Examples include recurring spending preferences or consistently
    expressed financial constraints.
    
    The distinction between explicit and implicit memory is motivated by
    prior work showing that real-world user interactions contain both direct
    and indirect personalization signals \cite{alpsbench}.
    
    \subsection{Memory Retrieval}
    
    Given a user query $q$, the memory subsystem retrieves a relevant subset:
    
    \begin{equation}
    R_q = R(M_u,q).
    \end{equation}
    
    The final LLM context is therefore represented as:
    
    \begin{equation}
    C = \{q,z,R_q\},
    \end{equation}
    
    where:
    
    \begin{itemize}
        \item $q$ is the user query,
        \item $z$ is deterministic financial state, and
        \item $R_q$ is retrieved personalized memory.
    \end{itemize}
    
    The LLM generates the final response:
    
    \begin{equation}
    y = LLM(C).
    \end{equation}
    
    \subsection{Model Context Protocol Interface}
    
    Switch FI exposes financial and memory operations through the Model
    Context Protocol (MCP).
    
    The MCP interface provides tools for operations such as:
    
    \begin{itemize}
        \item retrieving financial summaries,
        \item retrieving transactions,
        \item calculating category expenditure,
        \item querying monthly spending,
        \item retrieving user memories,
        \item storing explicit preferences,
        \item updating memory, and
        \item obtaining financial context for reasoning.
    \end{itemize}
    
    This makes the financial infrastructure independent of the specific
    LLM client.
    
    % =========================================================
    % 4. EVALUATION FRAMEWORK
    % =========================================================
    
    \section{Evaluation Framework}
    
    \subsection{Research Questions}
    
    The evaluation is organized around four research questions.
    
    \textbf{RQ1:} Does persistent memory improve preference adherence in
    personalized financial responses?
    
    \textbf{RQ2:} Does memory affect numerical grounding accuracy when
    financial facts are deterministically computed?
    
    \textbf{RQ3:} Does persistent memory improve retrieval of information
    expressed earlier in a conversation?
    
    \textbf{RQ4:} Does memory change the confidence or calibration behavior
    of the LLM when answering financial questions?
    
    \subsection{Dataset Construction}
    
    The evaluation uses a synthetic but realistic financial dataset generated
    through the Switch FI mock adapter.
    
    Each synthetic user is associated with one or more financial accounts
    containing transactions across multiple months.
    
    \begin{table}[t]
    \centering
    \caption{Evaluation Dataset Configuration}
    \label{tab:dataset}
    \begin{tabular}{lc}
    \toprule
    \textbf{Property} & \textbf{Value} \\
    \midrule
    Number of users & [TBD] \\
    Number of accounts & [TBD] \\
    Number of months & [TBD] \\
    Number of transactions & [TBD] \\
    Number of categories & [TBD] \\
    Evaluation questions & [TBD] \\
    Known ground truth & Yes \\
    \bottomrule
    \end{tabular}
    \end{table}
    
    The dataset is intentionally constructed so that every numerical
    quantity required by an evaluation question has a known ground-truth
    value.
    
    \subsection{Evaluation Configurations}
    
    The primary experiment is an ablation comparing two configurations.
    
    \begin{table}[t]
    \centering
    \caption{Experimental Configurations}
    \label{tab:ablation}
    \begin{tabular}{lll}
    \toprule
    \textbf{Configuration} & \textbf{Memory} & \textbf{Financial Tools} \\
    \midrule
    Memory-Off & Disabled & Enabled \\
    Memory-On & Enabled & Enabled \\
    \bottomrule
    \end{tabular}
    \end{table}
    
    Both configurations use:
    
    \begin{itemize}
        \item identical financial datasets,
        \item identical question sets,
        \item identical deterministic financial tools,
        \item identical LLM,
        \item identical system instructions.
    \end{itemize}
    
    The only manipulated variable is persistent memory availability.
    
    The evaluation model is:
    
    \begin{equation}
    \text{Claude Haiku 4.5 via OpenRouter}
    \end{equation}
    
    with the exact model identifier to be reported in the final experiment.
    
    \subsection{Evaluation Tasks}
    
    The question set contains approximately 30--50 fixed questions divided
    into three primary task categories.
    
    \begin{table*}[t]
    \centering
    \caption{Evaluation Task Categories}
    \label{tab:tasks}
    \begin{tabular}{p{0.22\textwidth}p{0.35\textwidth}p{0.30\textwidth}}
    \toprule
    \textbf{Task} & \textbf{Description} & \textbf{Example} \\
    \midrule
    
    Numerical Grounding &
    Tests whether the assistant returns the correct numerical value derived
    from financial records. &
    ``How much did I spend on food last month?'' \\
    
    Preference Adherence &
    Tests whether stored financial preferences influence recommendations. &
    ``How much should I allocate to discretionary spending given my savings
    goal?'' \\
    
    Memory Retrieval &
    Tests whether information introduced earlier is correctly recalled and
    used in a later interaction. &
    ``What savings goal did I tell you I was working toward?'' \\
    
    \bottomrule
    \end{tabular}
    \end{table*}
    
    \subsubsection{Task 1: Numerical Grounding}
    
    For a financial query $q$, the system produces an answer $\hat{a}$.
    The expected numerical answer is obtained from deterministic ground
    truth $a^*$.
    
    The numerical evaluation therefore measures:
    
    \begin{equation}
    Acc_{\mathrm{num}}
    =
    \frac{
    \sum_{i=1}^{N}
    \mathbb{1}[\hat{a}_i = a_i^*]
    }{
    N
    }.
    \end{equation}
    
    For values where rounding or formatting differences are expected, an
    absolute tolerance $\epsilon$ may be used:
    
    \begin{equation}
    \mathbb{1}
    [
    |\hat{a}_i-a_i^*| \leq \epsilon
    ].
    \end{equation}
    
    \subsubsection{Task 2: Preference Adherence}
    
    Preference adherence evaluates whether the response respects a known
    user preference or constraint.
    
    For question $q_i$ with expected preference behavior $p_i$, let
    
    \begin{equation}
    A_i =
    \begin{cases}
    1 & \text{if the response satisfies } p_i,\\
    0 & \text{otherwise}.
    \end{cases}
    \end{equation}
    
    Preference adherence rate is:
    
    \begin{equation}
    PAR =
    \frac{1}{N}
    \sum_{i=1}^{N} A_i.
    \end{equation}
    
    \subsubsection{Task 3: Memory Retrieval}
    
    For multi-turn questions, the assistant must retrieve information
    introduced in an earlier interaction.
    
    Let $m^*$ denote the ground-truth memory and $\hat{m}$ the information
    used by the model.
    
    Memory retrieval accuracy is:
    
    \begin{equation}
    Acc_{\mathrm{mem}}
    =
    \frac{
    \sum_{i=1}^{N}
    \mathbb{1}[\hat{m}_i = m_i^*]
    }{
    N
    }.
    \end{equation}
    
    Semantic matching may additionally be used where exact lexical matching
    is inappropriate.
    
    \subsection{Confidence Calibration}
    
    The assistant may additionally be asked to provide a confidence label
    for each answer.
    
    Let $c_i$ denote the predicted confidence and $r_i$ indicate whether the
    answer is correct.
    
    The relationship between confidence and correctness will be analyzed
    using:
    
    \begin{equation}
    Calibration =
    \frac{1}{N}
    \sum_{i=1}^{N}
    |c_i-r_i|.
    \end{equation}
    
    Lower values indicate closer agreement between stated confidence and
    observed correctness.
    
    % =========================================================
    % 5. EXPERIMENTAL RESULTS
    % =========================================================
    
    \section{Experimental Results}
    
    \subsection{Overall Results}
    
    The primary comparison between memory-enabled and memory-disabled
    configurations is shown in Table~\ref{tab:results}.
    
    \begin{table}[t]
    \centering
    \caption{Overall Evaluation Results}
    \label{tab:results}
    \begin{tabular}{lcc}
    \toprule
    \textbf{Metric} & \textbf{Memory-Off} & \textbf{Memory-On} \\
    \midrule
    Numerical Accuracy & [TBD] & [TBD] \\
    Preference Adherence & [TBD] & [TBD] \\
    Memory Retrieval & [TBD] & [TBD] \\
    Confidence Calibration & [TBD] & [TBD] \\
    \bottomrule
    \end{tabular}
    \end{table}
    
    \textbf{Important:} all values in this section remain placeholders until
    the experimental harness has been executed.
    
    \subsection{Numerical Grounding}
    
    The numerical grounding experiment measures whether introducing memory
    changes the model's ability to answer questions whose ground truth is
    derived from deterministic financial aggregation.
    
    \begin{table}[t]
    \centering
    \caption{Numerical Grounding Performance}
    \label{tab:numerical}
    \begin{tabular}{lccc}
    \toprule
    \textbf{Category} & \textbf{Off} & \textbf{On} & \textbf{$\Delta$} \\
    \midrule
    Income & [TBD] & [TBD] & [TBD] \\
    Expenses & [TBD] & [TBD] & [TBD] \\
    Categories & [TBD] & [TBD] & [TBD] \\
    Savings & [TBD] & [TBD] & [TBD] \\
    Overall & [TBD] & [TBD] & [TBD] \\
    \bottomrule
    \end{tabular}
    \end{table}
    
    \subsection{Preference Adherence}
    
    The second experiment evaluates whether persistent memory improves the
    assistant's ability to follow financial preferences.
    
    \begin{table}[t]
    \centering
    \caption{Preference Adherence Results}
    \label{tab:preference}
    \begin{tabular}{lccc}
    \toprule
    \textbf{Preference Type} & \textbf{Off} & \textbf{On} & \textbf{$\Delta$} \\
    \midrule
    Budget Constraint & [TBD] & [TBD] & [TBD] \\
    Savings Goal & [TBD] & [TBD] & [TBD] \\
    Spending Preference & [TBD] & [TBD] & [TBD] \\
    Risk Preference & [TBD] & [TBD] & [TBD] \\
    Overall & [TBD] & [TBD] & [TBD] \\
    \bottomrule
    \end{tabular}
    \end{table}
    
    \subsection{Memory Retrieval}
    
    Memory retrieval evaluates whether information introduced in an earlier
    turn is correctly used in a subsequent financial question.
    
    \begin{table}[t]
    \centering
    \caption{Multi-Turn Memory Retrieval}
    \label{tab:memory}
    \begin{tabular}{lccc}
    \toprule
    \textbf{Task} & \textbf{Off} & \textbf{On} & \textbf{$\Delta$} \\
    \midrule
    Explicit Memory & [TBD] & [TBD] & [TBD] \\
    Implicit Memory & [TBD] & [TBD] & [TBD] \\
    Multi-Hop Recall & [TBD] & [TBD] & [TBD] \\
    Overall & [TBD] & [TBD] & [TBD] \\
    \bottomrule
    \end{tabular}
    \end{table}
    
    \subsection{Ablation Analysis}
    
    The principal ablation compares the same model under two memory
    conditions.
    
    The relative improvement of memory-enabled reasoning is calculated as:
    
    \begin{equation}
    \Delta_{\mathrm{rel}}
    =
    \frac{
    Score_{\mathrm{on}}-Score_{\mathrm{off}}
    }{
    Score_{\mathrm{off}}
    }
    \times 100.
    \end{equation}
    
    \begin{figure}[t]
        \centering
        \includegraphics[
            width=\columnwidth
        ]{figures/evaluation_pipeline.pdf}
        \caption{
        Evaluation pipeline comparing memory-enabled and memory-disabled
        Switch FI configurations using identical financial data and question
        sets.
        }
        \label{fig:evaluation}
    \end{figure}
    
    % =========================================================
    % 6. DISCUSSION
    % =========================================================
    
    \section{Discussion}
    
    \subsection{Does Memory Improve Financial Personalization?}
    
    The primary purpose of the experiment is to determine whether persistent
    memory improves personalized financial reasoning.
    
    If the Memory-On configuration produces a higher preference adherence
    rate while maintaining comparable numerical accuracy, this would suggest
    that memory primarily improves personalization rather than financial
    grounding itself.
    
    \textbf{Observed result: [TBD after evaluation].}
    
    \subsection{Memory Versus Financial Grounding}
    
    The architecture deliberately separates memory from financial grounding.
    
    Deterministic financial tools provide:
    
    \begin{equation}
    \text{Ground Truth}_{financial}
    \end{equation}
    
    while memory provides:
    
    \begin{equation}
    \text{Context}_{personal}.
    \end{equation}
    
    The LLM combines these two sources during response generation.
    
    This distinction allows the experiment to determine whether memory
    introduces improvements independently of numerical grounding.
    
    \subsection{Comparison With Existing Personalization Benchmarks}
    
    The evaluation design is inspired by the observation in recent
    personalization benchmarks that remembering information and utilizing
    that information are separate capabilities \cite{alpsbench}.
    
    Unlike general conversational personalization benchmarks, Switch FI
    introduces a deterministic financial state as an additional source of
    ground truth.
    
    The resulting evaluation can therefore be viewed as:
    
    \begin{equation}
    \text{Personalization}
    +
    \text{Grounded Financial Reasoning}.
    \end{equation}
    
    This combination is the central distinction of the proposed framework.
    
    \subsection{Potential Failure Modes}
    
    Several failure modes are expected to be particularly relevant:
    
    \begin{enumerate}
    
        \item The model retrieves the correct memory but fails to apply it.
    
        \item The model applies a preference that conflicts with newer
        information.
    
        \item The model correctly retrieves financial data but performs
        incorrect reasoning over the values.
    
        \item The model generates unsupported numerical values despite
        receiving deterministic tool output.
    
        \item Implicit memory is inferred incorrectly from limited evidence.
    
    \end{enumerate}
    
    These failure modes should be analyzed separately rather than collapsing
    all errors into a single accuracy value.
    
    % =========================================================
    % 7. LIMITATIONS
    % =========================================================
    
    \section{Limitations}
    
    This study has several limitations.
    
    \subsection{Synthetic Evaluation Data}
    
    The experimental financial dataset is generated synthetically.
    Although the dataset is designed to resemble realistic financial
    transactions, it does not represent the complete variability of
    real-world banking data.
    
    \subsection{Sandbox Data Access}
    
    The current implementation uses the Finvu sandbox and a mock adapter
    rather than production financial accounts.
    
    Consequently, the evaluation does not establish the operational
    performance of the system against production-scale banking data.
    
    \subsection{Limited Model Coverage}
    
    The initial evaluation uses a single LLM configuration. Results may
    therefore depend on model-specific reasoning and instruction-following
    behavior.
    
    Future experiments should evaluate multiple models under identical
    conditions.
    
    \subsection{Evaluation Scale}
    
    The current evaluation uses a relatively small fixed question set.
    Although this provides controlled paired comparisons, a larger dataset
    would provide stronger statistical confidence.
    
    \subsection{Automated Evaluation}
    
    Some evaluation metrics rely on automated grading or deterministic
    matching. Human evaluation may be required for ambiguous cases involving
    preference interpretation and financial recommendations.
    
    % =========================================================
    % 8. CONCLUSION
    % =========================================================
    
    \section{Conclusion}
    
    This paper presented Switch FI, an MCP-based architecture for
    memory-augmented and financially grounded LLM assistance.
    
    The system separates financial truth from language generation by
    performing deterministic aggregation over structured transaction data,
    while a persistent memory layer provides user-specific preferences and
    context.
    
    We introduced an evaluation framework that compares memory-enabled and
    stateless configurations across numerical grounding, preference
    adherence, and multi-turn memory retrieval.
    
    The final experimental results will determine whether persistent memory
    provides measurable improvements in personalized financial reasoning
    without degrading numerical grounding.
    
    \textbf{Final empirical finding: [TBD].}
    
    Future work includes integration with production Account Aggregator
    credentials, evaluation across multiple LLMs, larger real-world user
    studies, improved implicit-memory handling, and deployment of Switch FI
    as a production financial MCP service.
    
    % =========================================================
    % ACKNOWLEDGMENTS
    % =========================================================
    
    \section*{Acknowledgment}
    
    The authors thank [Institution / Organization / Individuals] for
    supporting this work.
    
    % =========================================================
    % REFERENCES
    % =========================================================
    
    \bibliographystyle{IEEEtran}
    \bibliography{references}
    
    \end{document}